\documentclass[aspectratio=169,11pt]{beamer}
\makeatletter
\input{../../templates/dune-professional/beamerthemeDUNEProfessional.sty}
\makeatother
\title{VD detector protection}
\subtitle{Existing hardware, inhibit paths and VD additions}
\author{Manuel Arroyave}
\date{I\&I workshop draft --- 27 October 2026}
\newcommand{\mydate}{27 October 2026}
\hypersetup{colorlinks=true,urlcolor=DuneInfo,linkcolor=DuneInfo}
\newcommand{\planningnote}[1]{\par\medskip{\footnotesize\color{DuneTextMuted}#1\par}}
\begin{document}

\begin{frame}[plain]
  \titlepage
\end{frame}

\begin{frame}{DPS arbitrates protection; SC operates, monitors and logs}
  \renewcommand{\arraystretch}{1.5}
  \begin{tabularx}{\textwidth}{@{}>{\RaggedRight\arraybackslash}p{2.2cm}>{\RaggedRight\arraybackslash}X@{}}
    \textbf{DPS} & \textbf{Detector Protection System.} Uses facility and detector conditions to allow or block actions; issues permits, inhibits and trips.\\
    \textbf{SC} & \textbf{Slow controls, using Ignition.} Operates and monitors the detector; owns the data and event logs, including DPS decisions.\\
    \textbf{BMS} & \textbf{Building Management System.} Controls building services. FFDG BMS and \mbox{SDSTA BMS} are separate systems.\\
    \textbf{Local\newline interlocks} & \textbf{Equipment-level protection.} Can trip without waiting for SC or the central DPS.\\
  \end{tabularx}\par
  \bigskip
  \textbf{An operating request cannot override a protection trip.}
\end{frame}

\begin{frame}{DPS receives inputs from four system domains}
  \renewcommand{\arraystretch}{1.5}
  \begin{tabularx}{\textwidth}{@{}>{\RaggedRight\arraybackslash}p{2.7cm}>{\RaggedRight\arraybackslash}X@{}}
    \textbf{FFDG BMS} & API through the Batavia server at FNAL $\rightarrow$ SC $\rightarrow$ DPS.\\
    \textbf{SDSTA BMS} & Interface under investigation: available signals and connection points.\\
    \textbf{Cryogenics} & Selected process conditions. Each inhibit names the actual instrument or controller that supplies them.\\
    \textbf{Detector} & SC, the detector subsystems, and \textbf{rack infrastructure}: racks, PDUs and rack sensors.\\
  \end{tabularx}\par
  \bigskip
  Example: low liquid level withdraws the HV permit. SC logs the condition,
  DPS decision and equipment state.
  \planningnote{The FFDG route is the working integration basis. SDSTA retains personnel-safety responsibility. Facility and cryogenic systems retain control of their services.}
\end{frame}

\begin{frame}{VD uses SNMP control and dedicated protection I/O}
  \centering
  \begin{tikzpicture}[
    box/.style={fill=DunePaper,text=DuneOnAccent,align=center,
      text width=2.45cm,minimum height=1.35cm,inner sep=5pt,font=\small},
    edge/.style={-{Stealth},draw=DuneSky,line width=1.3pt},
    duplex/.style={arrows={Stealth-Stealth},draw=DuneSky,line width=1.3pt}]
    \node[box,text width=2.8cm,minimum height=1.0cm] (sc) at (-5.1,1.25) {\textbf{SC / Ignition}\\operations and logs};
    \node[box,fill=DuneSky,text width=2.8cm,minimum height=1.0cm] (plc) at (0,1.25) {\textbf{VD DPS PLC}\\CUC proposed};
    \node[box] (sub) at (-5.1,-1.35) {\textbf{Subsystems}\\OPC UA\\interfaces};
    \node[box] (io) at (-1.7,-1.35) {\textbf{BDE / bias I/O}\\input signals and\\inhibit outputs};
    \node[box] (lv) at (1.7,-1.35) {\textbf{Subsystem supplies}\\PL506 / ISEG\\PDS LV};
    \node[box,fill=DuneCoral] (rack) at (5.1,-1.35) {\textbf{Rack\\infrastructure}\\PDUs and\\rack sensors};
    \draw[duplex] (sc.east) -- node[above,font=\small] {OPC UA} (plc.west);
    \draw[duplex] (sc.south) -- node[right,font=\small] {OPC UA} (sub.north);
    \draw[duplex] (plc.south) -- (0,0.1) -- (-1.7,0.1) -- (io.north);
    \node[anchor=south,font=\small] at (-1.7,0.13) {EtherCAT};
    \draw[duplex] (plc.east) -- node[above,font=\small] {SNMP control / status} (5.1,1.25);
    \draw[draw=DuneSky,line width=1.3pt] (5.1,1.25) -- (5.1,0.1) -- (1.7,0.1);
    \draw[edge] (1.7,0.1) -- (lv.north);
    \draw[edge] (5.1,0.1) -- (rack.north);
    \draw[edge,draw=DuneCoral] (io.south) -- (-1.7,-2.5) --
      node[below,font=\small] {PL506 / ISEG wired inhibits} (1.7,-2.5) -- (lv.south);
  \end{tikzpicture}\par
  \medskip
  \DuneDiagramCaption{PDUs and WIENER supplies share SNMP control/status. BDE/bias also have dedicated wired protection.}
  \planningnote{PLC = programmable logic controller; CUC = Central Utility Cavern; LV = low voltage. Additional subsystem exchanges go through SC. TDE has PDU control now; other interfaces remain options. HD may extend this system or use another PLC.}
\end{frame}

\begin{frame}{BDE already defines a distributed protection system}
  BDE = bottom drift electronics. I/O = inputs and outputs.\par\medskip
  \centering
  \begin{tikzpicture}[
    box/.style={fill=DunePaper,text=DuneOnAccent,align=center,
      text width=3.1cm,minimum height=1.05cm,inner sep=6pt},
    edge/.style={-{Stealth},draw=DuneSky,line width=1.2pt}]
    \node[box,fill=DuneSky] (io) at (-4,0) {BDE auxiliary rack\\remote I/O};
    \node[box] (lv) at (1.6,1.1) {2 LV racks\\2 PL506 supplies};
    \node[box,fill=DuneCoral] (wiec) at (1.6,-1.1) {8 WIECs\\8 PTCs + 48 WIBs};
    \draw[edge] (io.north east) -| (-1.7,1.1) -- node[above,font=\small] {inhibits} (lv.west);
    \draw[edge,<->] (io.south east) -| (-1.7,-1.1) -- node[below,font=\small] {EtherCAT} (wiec.west);
    \draw[edge] (lv.east) -- (4.35,1.1) -- node[right,font=\small] {48 V} (4.35,-1.1) -- (wiec.east);
  \end{tikzpicture}\par
  \medskip
  \DuneDiagramCaption{Ten groups: 20 PL506 supplies, 80 WIECs and 480 WIBs.}
  \medskip
  \raggedright The auxiliary I/O also connects fans, heaters and temperature sensors.
  \planningnote{WIEC = warm interface electronics crate; PTC = power and timing card; WIB = warm interface board. June 2026 layout; fan/heater wiring omitted here.}
\end{frame}

\begin{frame}{BDE drawings, parts and software are already available}
  \begin{columns}[T,onlytextwidth]
    \begin{column}{0.47\textwidth}
      \textbf{TPC protection panel designs}\par\medskip
      \colorbox{DunePaper}{\includegraphics[page=2,width=\dimexpr\linewidth-2\fboxsep\relax,height=3.1cm,keepaspectratio,
        trim=290bp 500bp 900bp 430bp,clip]{figures/tpc-panel-reference.pdf}}\par
      \medskip\DuneDiagramCaption{Central PLC, bias I/O and LV/fan/heater reference panels.}
    \end{column}
    \begin{column}{0.49\textwidth}
      \textbf{BDE Ignition screens}\par\medskip
      \includegraphics[page=7,width=\linewidth,height=3.1cm,keepaspectratio,
        trim=10bp 210bp 460bp 74bp,clip]{figures/sc-progress-june2026.pdf}\par
      \planningnote{Existing display of WIEC groups; June 2026 SC progress presentation.}
    \end{column}
  \end{columns}\par
  \medskip
  Parts list: CX5140 PLC, EK1100 couplers, EK1521 fibre links and EL2124 inhibit outputs.
  \planningnote{Historical TPC reference: EDMS 2401090 v4, sheet 2, with parts and power workbooks. The VD baseline is eight WIECs per group.}
\end{frame}

\begin{frame}{Trace each inhibit to the hardware it uses}
  The matrix gives the condition and action. Trace the full path for each affected load.\par
  \bigskip
  \centering
  \begin{tikzpicture}[
    box/.style={fill=DunePaper,text=DuneOnAccent,align=center,
      text width=2.2cm,minimum height=1.35cm,inner sep=5pt,font=\small},
    edge/.style={-{Stealth},draw=DuneSky,line width=1.2pt}]
    \node[box] (source) at (-5.5,0) {Condition\\and source};
    \node[box] (input) at (-2.75,0) {Input signal\\and connection};
    \node[box,fill=DuneSky] (decision) at (0,0) {DPS / local\\decision};
    \node[box,fill=DuneCoral] (output) at (2.75,0) {Output\\and load};
    \node[box] (feedback) at (5.5,0) {Equipment state\\and feedback};
    \draw[edge] (source) -- (input);
    \draw[edge] (input) -- (decision);
    \draw[edge] (decision) -- (output);
    \draw[edge] (output) -- (feedback);
  \end{tikzpicture}\par
  \medskip
  \DuneDiagramCaption{Record the rack, connector, protocol, channel, power feed and allowed detection-to-shutdown time.}
  \bigskip
  \raggedright Use the existing rack lists, interfaces, panel drawings, parts and power calculations.\par
  \medskip
  \textbf{Result:} hardware already covered, additions with quantities, and open choices that change those quantities.
  \planningnote{Count shared inputs and outputs once, even when several inhibits use them. Wiring and software work are listed separately from hardware purchases.}
\end{frame}

\begin{frame}{BDE cooling already has an input-to-inhibit path}
  \centering
  \begin{tikzpicture}[
    box/.style={fill=DunePaper,text=DuneOnAccent,align=center,
      text width=1.75cm,minimum height=1.25cm,inner sep=4pt,font=\small},
    edge/.style={-{Stealth},draw=DuneSky,line width=1.2pt}]
    \node[box] (fan) at (-5.75,0.6) {Fan status\\DB25 cable};
    \node[box] (input) at (-3.45,0.6) {EL1004\\input I/O};
    \node[box,fill=DuneSky] (plc) at (-1.15,0.6) {DPS PLC\\cooling rule};
    \node[box] (output) at (1.15,0.6) {EL2124\\5 V output};
    \node[box,fill=DuneCoral] (supply) at (3.45,0.6) {PL506\\channel inhibit};
    \node[box] (load) at (5.75,0.6) {Affected WIEC\\48 V feed};
    \draw[edge] (fan) -- (input);
    \draw[edge] (input) -- (plc);
    \draw[edge] (plc) -- (output);
    \draw[edge] (output) -- (supply);
    \draw[edge] (supply) -- (load);
    \node[font=\small] at (-2.3,1.55) {EtherCAT};
    \node[font=\small] at (0,1.55) {EtherCAT};
    \node[font=\small] at (2.3,1.55) {Wired};
  \end{tikzpicture}\par
  \medskip
  \DuneDiagramCaption{Cooling fault $\rightarrow$ withdraw the permit for the affected PL506 channel.}
  \medskip
  \raggedright \textbf{Feedback:} PL506 state and voltage/current $\rightarrow$ SNMP $\rightarrow$ DPS;
  OPC UA carries the result to SC for logging.\par
  \medskip
  \textbf{VD allocation check:} six fan-status signals per WIEC means 48 per
  eight-WIEC group. Match these to the populated VD inputs and cables.\par
  \planningnote{The six-fan WIEC note differs from the four-fan ICEBERG reference. These are baseline requirements, not new-module counts. Fan voting and allowed shutdown time remain to be agreed.}
\end{frame}

\begin{frame}{Rack infrastructure is a separate protection system}
  \centering
  \begin{tikzpicture}[
    box/.style={fill=DunePaper,text=DuneOnAccent,align=center,
      text width=2.0cm,minimum height=1.05cm,inner sep=5pt,font=\small},
    edge/.style={-{Stealth},draw=DuneSky,line width=1.2pt},
    duplex/.style={arrows={Stealth-Stealth},draw=DuneSky,line width=1.2pt}]
    \draw[draw=DuneBorder,line width=1pt,fill=DuneSurface] (0.3,-0.9) rectangle (3.7,2.45);
    \node[font=\small\bfseries] at (2,2.2) {Rack infrastructure};
    \node[box,text width=2.2cm] (sc) at (-5.5,0) {SC / Ignition\\operations\\and logs};
    \node[box,fill=DuneSky] (plc) at (-1.8,0) {DPS PLC\\arbitration};
    \node[box,fill=DuneCoral,text width=2.4cm] (pdu) at (2,0) {PDU\\local trips\\outlet control};
    \node[box,text width=1.8cm] (load) at (5.6,0) {Subsystem\\loads};
    \node[box,text width=2.8cm,minimum height=0.9cm] (sensors) at (2,1.45) {Smoke + humidity\\temperature / leak};
    \draw[duplex] (sc.east) -- node[above,font=\small] {OPC UA} (plc.west);
    \draw[duplex] (plc.east) -- node[above,font=\small] {SNMP} (pdu.west);
    \draw[edge] (pdu.east) -- node[above,font=\small] {AC power} (load.west);
    \draw[edge] (sensors.south) -- (pdu.north);
  \end{tikzpicture}\par
  \medskip
  \DuneDiagramCaption{Count the rack, PDUs, sensors and interface parts once. Link each subsystem to its PDU outlets.}
  \medskip
  \raggedright Local PDU rules protect the mapped outlets. Humidity is monitored;
  a humidity shutdown needs a stated requirement.\par
  \planningnote{The PX3 package supplies the smoke interface, powered contact hub and relay. Temperature/leak sensors are included where fitted. Reset and re-enable remain separate actions.}
\end{frame}

\begin{frame}{Detector additions follow the existing VD hardware}
  \small
  \renewcommand{\arraystretch}{1.2}
  \begin{tabularx}{\textwidth}{@{}>{\RaggedRight\arraybackslash}p{1.9cm}>{\RaggedRight\arraybackslash}p{4.7cm}>{\RaggedRight\arraybackslash}X@{}}
    \toprule
    \textbf{System} & \textbf{Starting hardware} & \textbf{What determines additions}\\
    \midrule
    BDE / HVS & BDE I/O; PL506 / ISEG inhibits; Heinzinger and cRIO & Signal and load mapping against the existing VD design\\
    PDS / PoF & 34 LV feeds; 80 PoF boxes & Opta PLC or Beckhoff PLC: local I/O, power, enclosures and cables\\
    TDE & 20 racks; current control through PDUs & Crate status / external inhibit retained as future options\\
    CalCI & Existing instrumentation and calibration descriptions & Status, optical permits and actuator connections\\
    \bottomrule
  \end{tabularx}\par
  \medskip
  \normalsize The 80 PoF boxes have 1,280 monitoring values. Each controller
  option must cover the same physical signals and required actions.
  \planningnote{PDS = photon detection system; PoF = power over fibre; TDE = top drift electronics; HVS = high-voltage system; CalCI = calibration and cryogenic instrumentation. These values are not a central DPS terminal count.}
\end{frame}

\begin{frame}{Facility interfaces need named signals and routes}
  \renewcommand{\arraystretch}{1.4}
  \begin{tabularx}{\textwidth}{@{}>{\RaggedRight\arraybackslash}p{2.2cm}>{\RaggedRight\arraybackslash}X@{}}
    \toprule
    \textbf{Source} & \textbf{Hardware consequence of the path review}\\
    \midrule
    FFDG BMS & Use the API through Batavia and SC. Name the required points; the API exchange itself adds no discrete or analog DPS terminals.\\
    SDSTA BMS & The survey identifies available signals, controllers and connection points. Those determine any interface hardware.\\
    Cryogenics & Name the measurement or controller status and its transport. Some level measurements come through detector instrumentation and SC.\\
    \bottomrule
  \end{tabularx}\par
  \medskip
  \textbf{FAC-01 example:} map the fire-alarm signal to the drift-HV and CRP/FC
  bias shutdown paths before assigning an additional facility input.
  \planningnote{The known FFDG API route does not establish that every required alarm is exposed on it. SDSTA interface work remains under investigation.}
\end{frame}

\begin{frame}{Test arbitration, trips and recovery on the bench}
  \centering
  \begin{tikzpicture}[
    box/.style={fill=DunePaper,text=DuneOnAccent,align=center,
      text width=2.25cm,minimum height=1.0cm,inner sep=5pt},
    edge/.style={-{Stealth},draw=DuneSky,line width=1.2pt}]
    \node[box] (in) at (-4.8,0) {Fault input or\\conflicting request};
    \node[box,fill=DuneSky] (plc) at (-1.6,0) {DPS PLC\\arbitration};
    \node[box,fill=DuneCoral] (out) at (1.6,0) {Supply inhibit\\or PDU command};
    \node[box] (state) at (4.8,0) {Equipment state\\and feedback};
    \draw[edge] (in) -- (plc);
    \draw[edge] (plc) -- (out);
    \draw[edge] (out) -- (state);
  \end{tikzpicture}\par
  \medskip
  \DuneDiagramCaption{Use the BDE vertical-slice work and G271 equipment for representative VD paths.}
  \bigskip
  Measure detection through completed shutdown. Test local trips, SC and fieldbus
  loss, reset and re-enable. Verify that SC logs the decisions and equipment states.\par
  \bigskip
  No critical sub-ms function has been identified so far. Simulation prepares the
  fault and load cases; hardware tests measure the complete response.
\end{frame}

\begin{frame}{Take the VD hardware to review in June 2027}
  \renewcommand{\arraystretch}{1.4}
  \begin{tabularx}{\textwidth}{@{}>{\RaggedRight\arraybackslash}p{3.2cm}>{\RaggedRight\arraybackslash}X@{}}
    \toprule
    \textbf{Work before review} & \textbf{Result for procurement}\\
    \midrule
    Trace the VD inhibits & Existing hardware, added channels and cables, quantities and open alternatives\\
    Build representative paths & Measured trip time, failure/recovery results and SC logging\\
    Prepare purchases & Panel space and power, drawings, parts lists, costs, lead times and spares\\
    \bottomrule
  \end{tabularx}\par
  \bigskip
  After readiness review: order the hardware and install it in the site sequence.
  \planningnote{Possible facility-interface need: summer 2027. Full VD response-matrix review: early 2028. The current path review establishes the requirements that affect hardware purchases.}
\end{frame}

\begin{frame}{I will lead the VD hardware work with a control engineer}
  \renewcommand{\arraystretch}{1.4}
  \begin{tabularx}{\textwidth}{@{}>{\RaggedRight\arraybackslash}p{4.1cm}>{\RaggedRight\arraybackslash}X@{}}
    \toprule
    \textbf{Proposed contribution} & \textbf{Work}\\
    \midrule
    Manuel + control engineer & VD design, SURF interfaces and procurement package\\
    Terri and Trevor & BDE panels, PLC/I/O implementation and bench work\\
    Geoff + Manuel & SC commands, PDU protection and recovery\\
    Vivian + consortia / DPS & Equipment requirements and ICD updates\\
    \bottomrule
  \end{tabularx}\par
  \bigskip
  \textbf{Support needed:} control-engineering time and access to the existing test stands.
  \planningnote{SURF work: inspection maps are ready; access is being coordinated with Rachel, with Dale Curran identified as an SDSTA integration contact.}
\end{frame}

\appendix
\begin{frame}{Backup: VD rack space and the existing DPS allocation}
  \renewcommand{\arraystretch}{1.4}
  \begin{tabularx}{\textwidth}{@{}>{\RaggedRight\arraybackslash}p{3.4cm}r>{\RaggedRight\arraybackslash}X@{}}
    \toprule
    \textbf{Rack type} & \textbf{Count} & \textbf{DPS allocation and workbook space}\\
    \midrule
    BDE auxiliary & 10 & 12U DPS; U3--U21 marked Space (19U); eight WIECs per rack\\
    BDE low voltage & 20 & One PL506 per rack; four WIECs per supply\\
    Bias high voltage & 10 & 12U DPS; U3--U12 marked Space (10U)\\
    \bottomrule
  \end{tabularx}\par
  \bigskip
  Space marked in the rack workbook is not a new DPS reservation.
  Free DIN-rail length and I/O capacity depend on the populated VD panel.
  \planningnote{May 2025 VD workbook: BDE Aux Rack and Bias HV Rack. Auxiliary U9 also has a possible ground-plane readout note. U is rack height.}
\end{frame}

\begin{frame}{The VD bias-panel draft specifies two DIN rails per rack}
  \renewcommand{\arraystretch}{1.3}
  \begin{tabularx}{\textwidth}{@{}>{\RaggedRight\arraybackslash}Xrr@{}}
    \toprule
    \textbf{Component} & \textbf{Per rail} & \textbf{10 racks}\\
    \midrule
    EK1100 EtherCAT coupler & 1 & 20\\
    EL9410 supply terminal & 1 & 20\\
    EL2124: four 5 V outputs & 18 & 360\\
    EL9011 end cap & 1 & 20\\
    24 V power supply & 1 & 20\\
    \bottomrule
  \end{tabularx}\par
  \medskip
  Each rail occupies 337 mm. The draft provides 1,440 output channels;
  the working VD inventory lists 576 bias channels.\par
  \planningnote{May 2025 workbook, Bias HV Rack!K17:K21. This is the first-pass layout; the difference from 576 bias channels is not verified spare capacity.}
\end{frame}

\begin{frame}{Use the component loads for the VD power budget}
  \small
  \renewcommand{\arraystretch}{1.3}
  \begin{tabularx}{\textwidth}{@{}>{\RaggedRight\arraybackslash}p{3.7cm}>{\RaggedRight\arraybackslash}p{2.4cm}>{\RaggedRight\arraybackslash}X@{}}
    \toprule
    \textbf{Equipment} & \textbf{Power / capacity} & \textbf{Basis and limit}\\
    \midrule
    PL506 input & 2.25 kW each & 20 supplies: about 45.1 kW expected AC input\\
    WIEC fans only & 380 W per group & $8\times6\times0.33\,\mathrm{A}\times24\,\mathrm{V}$; excludes heaters and I/O\\
    Bias DPS supply allowance & 24 V, 40 A & 960 W DC capacity per rack in the workbook\\
    \bottomrule
  \end{tabularx}\par
  \medskip
  The inherited 831 W auxiliary figure needs a VD component budget
  covering fans, heaters and populated I/O.\par
  \planningnote{May workbook; WIEC safety note, internal revision 0. Fan arithmetic assumes six fans per WIEC. Rated load, expected AC input and DC supply capacity are distinct.}
\end{frame}

\begin{frame}{Rack infrastructure already has a hardware reference}
  \small
  \renewcommand{\arraystretch}{1.2}
  \begin{tabularx}{\textwidth}{@{}>{\RaggedRight\arraybackslash}p{3.4cm}>{\RaggedRight\arraybackslash}X@{}}
    \toprule
    \textbf{Rack component} & \textbf{Reference hardware / allocation}\\
    \midrule
    PDU & Raritan PX3; inlet and outlet allocation for the connected loads\\
    Smoke detector & 4WTAR-B and EOLR-1 end-of-line relay\\
    Contact interface & DX-PD2C5; 12 V supply to the smoke detector\\
    Humidity & Record the probe and peripheral interface for each rack\\
    Temperature / water leak & PDU peripheral sensors where fitted\\
    \bottomrule
  \end{tabularx}\par
  \medskip
  \normalsize The May workbook includes a 76-entry rack-protection/PDU allowance.
  Map the actual rack packages against this existing allowance.\par
  \planningnote{DUNE-doc-32655 v8, sections 3 and 5; May workbook, Standard Rack and Power Summary!B28. The guide does not specify a humidity probe model or trip requirement.}
\end{frame}

\begin{frame}{Reference: local PDU trips and recovery}
  \renewcommand{\arraystretch}{1.45}
  \begin{tabularx}{\textwidth}{@{}>{\RaggedRight\arraybackslash}p{3.7cm}>{\RaggedRight\arraybackslash}X@{}}
    \toprule
    \textbf{PX3 setup guide} & \textbf{Local behaviour documented}\\
    \midrule
    Smoke, water, temperature & Local event rules and scripts can switch off the protected outlets\\
    Network-switch outlet & Kept powered during the guide's fallback shutdown\\
    After a local trip & Power is restored deliberately; no automatic restoration\\
    \bottomrule
  \end{tabularx}\par
  \bigskip
  The VD control chain is Ignition $\rightarrow$ DPS PLC $\rightarrow$ PDU.
  Local trip status returns to the PLC and to SC for logging.
  \planningnote{DUNE-doc-32655 v8 describes an Ignition-direct prototype with local fallback on loss of Ignition. The VD PLC command path requires matching reachability, trip and reset configuration; it is not a claim that the prototype already implements it.}
\end{frame}

\begin{frame}{The VD inventory covers the other subsystem connections}
  \renewcommand{\arraystretch}{1.35}
  \begin{tabularx}{\textwidth}{@{}>{\RaggedRight\arraybackslash}p{3.2cm}>{\RaggedRight\arraybackslash}X@{}}
    \toprule
    \textbf{Subsystem} & \textbf{Working inventory and interfaces}\\
    \midrule
    PDS readout / calibration & 34 LV feeds in the current presentation scope; supply status and control interfaces\\
    Power over fibre (PoF) & 80 boxes; 640 temperature and 640 voltage readings; Opta PLC or Beckhoff PLC alternative\\
    TDE & 20 mezzanine racks; PDU control now; other TDE interfaces are possible extensions\\
    Calibration / cryogenic instrumentation & Level and temperature readout, purity monitors, ionization laser and neutron source\\
    \bottomrule
  \end{tabularx}
  \planningnote{Additional subsystem exchanges use OPC UA through SC. June 2026 SC Progress, slide 29, names the PoF controller alternatives. Box count and monitoring values do not establish controller or central DPS terminal quantities.}
\end{frame}

\begin{frame}{The BMS drawing identifies 20 control-panel symbols}
  \renewcommand{\arraystretch}{1.4}
  \begin{tabularx}{\textwidth}{@{}>{\RaggedRight\arraybackslash}p{3.4cm}>{\RaggedRight\arraybackslash}X@{}}
    \toprule
    \textbf{Drawing area} & \textbf{Interfaces on the inspection maps}\\
    \midrule
    CUC, level 4850 & Eight panel references: air handling, chilled water, power and associated equipment\\
    Level 4910 & North and south sump interfaces, SP-1 and SP-2\\
    4850-08-ES group & Spare panel for future equipment; mapped to the access-drift inspection area\\
    Other 4850 areas & Chillers, pumps, ventilation and generator interfaces\\
    \bottomrule
  \end{tabularx}\par
  \medskip
  Inspection planned with Rachel; Dale Curran is the SDSTA integration contact identified for this work.
  \planningnote{U-FD-M-901 rev12, 14 February 2025; marked 4850/4910 maps and survey register. Symbols are drawing references; field observations are pending.}
\end{frame}

\begin{frame}{Reference: the May 2025 VD rack inventory}
  \renewcommand{\arraystretch}{1.15}
  \begin{tabularx}{\textwidth}{@{}>{\RaggedRight\arraybackslash}Xr>{\RaggedRight\arraybackslash}Xr@{}}
    \toprule
    \textbf{Rack type} & \textbf{Count} & \textbf{Rack type} & \textbf{Count}\\
    \midrule
    BDE low voltage & 20 & PDS low voltage & 4\\
    BDE auxiliary & 10 & Drift HV & 2\\
    TDE & 20 & Reserved drift HV & 2\\
    Bias HV & 10 & Spare / CalCI & 8\\
    \bottomrule
  \end{tabularx}\par
  \bigskip
  \textbf{76 numbered mezzanine positions:} SEDM-01 to SEDM-76.\par
  \medskip
  The worksheet also has a separate concept for 40 top-plate mini-racks.
  \planningnote{May 2025 allocation: rack positions, distinct from PLC panels and remote-I/O stations. Includes the listed reserved and spare positions.}
\end{frame}

\end{document}
